% ---------------------------------------------------------------
%  Preámbulo — configuración general del documento
% ---------------------------------------------------------------

% ---- Idioma y codificación ------------------------------------
\usepackage[utf8]{inputenc}
\usepackage[T1]{fontenc}
\usepackage[spanish,es-tabla,es-nodecimaldot]{babel}
\usepackage{csquotes}

% ---- Tipografía -----------------------------------------------
% El documento original usa Calibri. "carlito" es la fuente libre
% métricamente compatible; si no está instalada se usa la de por defecto.
\IfFileExists{carlito.sty}{\usepackage{carlito}\renewcommand{\familydefault}{\sfdefault}}{}
\usepackage[protrusion=true,expansion=false]{microtype}
\usepackage{amssymb}      % casillas del cuestionario ($\square$)

% ---- Página y espaciado (idénticos al .docx) -------------------
% A4, márgenes: sup. 2,5 cm · inf. 2,5 cm · izq. 3 cm · der. 2 cm
\usepackage[a4paper,top=2.5cm,bottom=2.5cm,left=3cm,right=2cm]{geometry}
\usepackage{setspace}
\onehalfspacing                       % interlineado 1,5
\setlength{\parindent}{0pt}
\setlength{\parskip}{6pt}

% ---- Color corporativo UNIR ------------------------------------
\usepackage{xcolor}
\definecolor{unirazul}{HTML}{0098CD}

% ---- Títulos numerados y en azul --------------------------------
\usepackage{titlesec}
\titleformat{\section}{\normalfont\Large\bfseries\color{unirazul}}{\thesection.}{0.5em}{}
\titleformat{\subsection}{\normalfont\large\bfseries\color{unirazul}}{\thesubsection.}{0.5em}{}
\titleformat{\subsubsection}{\normalfont\normalsize\bfseries\color{unirazul}}{\thesubsubsection.}{0.5em}{}
\titlespacing*{\section}{0pt}{18pt}{10pt}
\titlespacing*{\subsection}{0pt}{14pt}{8pt}
\titlespacing*{\subsubsection}{0pt}{12pt}{6pt}

% ---- Índice ----------------------------------------------------
\usepackage{tocloft}
\setlength{\cftsecnumwidth}{2.2em}   % hueco del número de sección
\setcounter{tocdepth}{3}
\setcounter{secnumdepth}{3}

% ---- Figuras, tablas y listas ----------------------------------
\usepackage{graphicx}
\graphicspath{{imagenes/}}
\usepackage{booktabs}
\usepackage{tabularx}
\usepackage{longtable}
\usepackage{array}
\usepackage{calc}          % anchos de columna calculados (tablas de pandoc)
\usepackage{float}
\usepackage{multirow}
\usepackage[font=small,labelfont=bf]{caption}
\usepackage{enumitem}

% Pie de tabla/figura fuera de flotante y línea de fuente
\captionsetup{justification=raggedright,singlelinecheck=false}
\newcommand{\fuente}[1]{{\par\raggedright\footnotesize\itshape #1\par}\medskip}
\setlist{itemsep=2pt,topsep=4pt}

% ---- Encabezado y pie ------------------------------------------
\usepackage{fancyhdr}
\pagestyle{fancy}
\fancyhf{}
\fancyfoot[L]{\footnotesize\color{gray}\autorescorto}
\fancyfoot[R]{\footnotesize\thepage}
\renewcommand{\headrulewidth}{0pt}
\renewcommand{\footrulewidth}{0.4pt}

% ---------------------------------------------------------------
%  BIBLIOGRAFÍA
%  Deja \overleaftrue para compilar en Overleaf con estilo APA 7
%  (biblatex-apa + biber). Cambia a \overleaffalse solo si tu
%  instalación local no tiene biblatex (usa natbib + apalike).
% ---------------------------------------------------------------
\newif\ifoverleaf
\overleaftrue

\ifoverleaf
  \usepackage[style=apa,natbib=true,backend=biber,language=spanish]{biblatex}
  \DeclareLanguageMapping{spanish}{spanish-apa}
  \addbibresource{bibliografia/referencias.bib}
  \newcommand{\imprimirbibliografia}{%
    \printbibliography[title={Referencias bibliográficas},heading=bibnumbered]}
\else
  \usepackage{natbib}
  \bibliographystyle{apalike}
  \setcitestyle{authoryear,round}
  \newcommand{\imprimirbibliografia}{%
    \renewcommand{\bibname}{Referencias bibliográficas}%
    \bibliography{bibliografia/referencias}}
\fi

% ---- Hipervínculos (siempre al final del preámbulo) -------------
\usepackage[hidelinks,breaklinks=true]{hyperref}
\usepackage{url}
\urlstyle{same}
